%+TITLE: Teorías de gauge no abelianas
%+DESCRIPTION: Construimos teorías de Yang--Mills y describimos la interacción fuerte mediante la simetría de gauge $SU(3)$
% Fuente: pag-31-35/transcripcion_teoria_clasica_campos_paginas_31-35/transcripcion_paginas_31-35.tex
% Páginas manuscritas 34--35.
% Se conserva la transcripción; los puntos pendientes de revisión se anotan abajo.
% Véase también pag-31-35/transcripcion_teoria_clasica_campos_paginas_31-35/NOTAS_TRANSCRIPCION.md.
% Página 34, factor \(g\) en las rotaciones: el manuscrito parametriza \(R=e^{ig\vec\theta\cdot\vec J}\). Se conservó ese factor aun cuando en una discusión puramente de \(SO(3)\) suele absorberse en los parámetros.
% Página 35, generadores de \(SU(3)\): se reprodujeron las ocho matrices \(\lambda^{(i)}\) exactamente como matrices de Gell--Mann, sin introducir el factor habitual \(1/2\) que a veces se usa al definir los generadores.
% Página 35, frase final: se conservó «objetos escalares \(SU(3)\)» tal como está escrito; probablemente convenga decidir al editar si se quería decir singletes de \(SU(3)\).
\documentclass[11pt,a4paper]{article}

\usepackage[utf8]{inputenc}
\usepackage[T1]{fontenc}
\usepackage[spanish,es-nodecimaldot]{babel}
\usepackage{lmodern}
\usepackage{amsmath,amssymb,mathtools}
\usepackage{graphicx}
\usepackage{xcolor}
\usepackage{geometry}
\usepackage{enumitem}
\usepackage{microtype}

\geometry{margin=2.3cm}
\setlength{\parindent}{0pt}
\setlength{\parskip}{0.65em}
\setlist[itemize]{leftmargin=1.6em,itemsep=0.25em,topsep=0.25em}

\definecolor{title-red}{RGB}{165,25,20}
\definecolor{concept-green}{RGB}{20,125,45}
\newcommand{\concept}[1]{\textcolor{concept-green}{#1}}

\newcommand{\ket}[1]{\left|#1\right\rangle}
\newcommand{\bra}[1]{\left\langle#1\right|}
\newcommand{\braket}[2]{\left\langle #1\middle|#2\right\rangle}
\newcommand{\hc}{\mathrm{h.c.}}
\newcommand{\Tr}{\operatorname{Tr}}
\newcommand{\Diffs}{\mathrm{Diffs}}
\newcommand{\slashedpartial}{\not\!\partial}

\begin{document}

% -----------------------------------------------------------------------------
% Página manuscrita 34
% -----------------------------------------------------------------------------
Inspirados en esta manera de entender la electrodinámica, a mediados del siglo
XX describieron las interacciones nucleares usando \concept{simetrías de gauge}.

Hagamos un ejemplo con campos escalares. Supongamos que tenemos tres campos
$\phi_1,\phi_2,\phi_3$ y el lagrangiano
\[
  \mathcal L=
  -\frac12\partial_\mu\phi_1\partial^\mu\phi_1
  -\frac12\partial_\mu\phi_2\partial^\mu\phi_2
  -\frac12\partial_\mu\phi_3\partial^\mu\phi_3
  -\frac{m^2}{2}\phi_1^2
  -\frac{m^2}{2}\phi_2^2
  -\frac{m^2}{2}\phi_3^2.
\]
Lo podemos escribir de forma más compacta si definimos
\[
  \vec\Phi=
  \begin{pmatrix}
    \phi_1\\
    \phi_2\\
    \phi_3
  \end{pmatrix},
  \qquad
  \mathcal L=-\frac12\partial_\mu\vec\Phi\cdot\partial^\mu\vec\Phi
             -\frac{m^2}{2}\vec\Phi\cdot\vec\Phi.
\]

Vemos que este lagrangiano tiene una simetría $SO(3)$,
\[
  \vec\Phi\longmapsto R\vec\Phi.
\]
Por ejemplo,
\[
  \partial_\mu\vec\Phi^{\,T}\partial^\mu\vec\Phi'
  =\partial_\mu\vec\Phi^{\,T}R^TR\partial^\mu\vec\Phi
  =\partial_\mu\vec\Phi^{\,T}\partial^\mu\vec\Phi,
\]
etcétera.

Recordemos que
\[
  R=e^{ig\vec\theta\cdot\vec J},
  \qquad
  \vec\theta=\theta\hat n=(\theta_1,\theta_2,\theta_3).
\]
Análogo al caso de electrodinámica, queremos que los parámetros $\theta_i$
puedan depender del espacio-tiempo. El lagrangiano no es invariante en este
caso:
\[
  \partial_\mu\vec\Phi'
  =\partial_\mu\left(e^{ig\vec\theta\cdot\vec J}\vec\Phi\right)
  =\partial_\mu\left(e^{ig\vec\theta\cdot\vec J}\right)\vec\Phi
   +e^{ig\vec\theta\cdot\vec J}\partial_\mu\vec\Phi.
\]
El nuevo término hace que
$\partial_\mu\vec\Phi\cdot\partial^\mu\vec\Phi$ no sea invariante. Para
construir objetos invariantes necesitamos una \concept{derivada covariante}
tal que
\[
  D_\mu'\vec\Phi'
  =e^{ig\vec\theta\cdot\vec J}D_\mu\vec\Phi.
\]
Definimos
\[
  D_\mu=\partial_\mu-ig\,\vec J\cdot\vec W_\mu,
\]
introduciendo tres bosones de gauge, y tomamos
\[
  \vec W_\mu'\cdot\vec J
  =e^{ig\vec\theta\cdot\vec J}
     (\vec W_\mu\cdot\vec J)
     e^{-ig\vec\theta\cdot\vec J}
  -\frac{i}{g}
   \partial_\mu\left(e^{ig\vec\theta\cdot\vec J}\right)
   e^{-ig\vec\theta\cdot\vec J}.
\]
Entonces
\begin{align*}
  D_\mu'\vec\Phi'
  &=(\partial_\mu-ig\vec J\cdot\vec W_\mu')
    \left(e^{ig\vec\theta\cdot\vec J}\vec\Phi\right)\\
  &=e^{ig\vec\theta\cdot\vec J}\partial_\mu\vec\Phi
    +\partial_\mu\left(e^{ig\vec\theta\cdot\vec J}\right)\vec\Phi
    -ig e^{ig\vec\theta\cdot\vec J}
       (\vec W_\mu\cdot\vec J)\vec\Phi
    -\partial_\mu\left(e^{ig\vec\theta\cdot\vec J}\right)\vec\Phi\\
  &=e^{ig\vec\theta\cdot\vec J}
    \left(\partial_\mu-ig\vec W_\mu\cdot\vec J\right)\vec\Phi
   =e^{ig\vec\theta\cdot\vec J}D_\mu\vec\Phi.
\end{align*}

Reemplazando las derivadas por derivadas covariantes tenemos un lagrangiano
invariante, pero introdujimos tres bosones de gauge $\vec W_\mu$. Necesitamos
un término cinético para ellos. Para construirlo, notemos que podemos
interpretar la derivada covariante como si transformara de la siguiente manera:
\[
  D_\mu\longmapsto D_\mu'
  =e^{ig\vec\theta\cdot\vec J}D_\mu e^{-ig\vec\theta\cdot\vec J}.
\]
Entonces el conmutador de derivadas transforma como
\[
  [D_\mu,D_\nu]\longmapsto[D_\mu',D_\nu']
  =e^{ig\vec\theta\cdot\vec J}
    [D_\mu,D_\nu]
    e^{-ig\vec\theta\cdot\vec J}.
\]
Es más, este conmutador no depende del campo sobre el cual actúa.

% -----------------------------------------------------------------------------
% Página manuscrita 35
% -----------------------------------------------------------------------------
\begin{align*}
  [D_\mu,D_\nu]\vec\Phi
  ={}&
  (\partial_\mu-ig\vec J\cdot\vec W_\mu)
  (\partial_\nu-ig\vec J\cdot\vec W_\nu)\vec\Phi\\
  &-(\partial_\nu-ig\vec J\cdot\vec W_\nu)
  (\partial_\mu-ig\vec J\cdot\vec W_\mu)\vec\Phi\\[0.3em]
  ={}&
  \partial_\mu\partial_\nu\vec\Phi
  -ig\vec J\cdot(\partial_\mu\vec W_\nu)\vec\Phi
  -ig\vec J\cdot\vec W_\nu\partial_\mu\vec\Phi
  -ig\vec J\cdot\vec W_\mu\partial_\nu\vec\Phi
  -g^2(\vec J\cdot\vec W_\mu)(\vec J\cdot\vec W_\nu)\vec\Phi\\
  &-\partial_\nu\partial_\mu\vec\Phi
  +ig\vec J\cdot(\partial_\nu\vec W_\mu)\vec\Phi
  +ig\vec J\cdot\vec W_\mu\partial_\nu\vec\Phi
  +ig\vec J\cdot\vec W_\nu\partial_\mu\vec\Phi
  +g^2(\vec J\cdot\vec W_\nu)(\vec J\cdot\vec W_\mu)\vec\Phi\\
  ={}&-ig\left[
       \vec J\cdot(\partial_\mu\vec W_\nu-\partial_\nu\vec W_\mu)
       -igW_\mu^aW_\nu^b[J^a,J^b]
     \right]\vec\Phi\\
  ={}&-igJ^c\left(
      \partial_\mu W_\nu^c-\partial_\nu W_\mu^c
      +g\epsilon_{abc}W_\mu^aW_\nu^b
    \right)\vec\Phi\\
  \equiv{}&-igJ^cF_{\mu\nu}^c\vec\Phi.
\end{align*}
El objeto $F_{\mu\nu}^c$ es el \concept{tensor de fuerza del campo}.

De la discusión anterior, este tensor de fuerza transforma como
\[
  \vec F_{\mu\nu}'\cdot\vec J
  =e^{ig\vec\theta\cdot\vec J}
    (\vec F_{\mu\nu}\cdot\vec J)
    e^{-ig\vec\theta\cdot\vec J}.
\]
Esto quiere decir que si $F_{\mu\nu}^c=0$ en un gauge, será cero también en otro
gauge. Para obtener un objeto invariante de Lorentz y de gauge sólo tenemos que
hacer
\[
  -\frac14\Tr\left[
    (\vec F_{\mu\nu}\cdot\vec J)
    (\vec F^{\mu\nu}\cdot\vec J)
  \right]
  =-\frac14\Tr(F_{\mu\nu}F^{\mu\nu}),
  \qquad
  F_{\mu\nu}\equiv\vec J\cdot\vec F_{\mu\nu}.
\]

A diferencia del caso $U(1)$, ahora los bosones de gauge tienen
auto-interacciones; es decir, el lagrangiano contiene términos con tres o
cuatro $W_\mu$'s. Estas teorías son no lineales. Estas teorías se llaman
teorías de \concept{Yang--Mills}.

Una de las teorías de Yang--Mills más estudiadas describe la interacción
nuclear fuerte entre quarks. Es una teoría de gauge $SU(3)$. El quark $u$ se
escribe como un triplete de espinores de Dirac
\[
  \begin{pmatrix}
    u_1\\
    u_2\\
    u_3
  \end{pmatrix}
\]
con carga electromagnética $q=\frac23e$. Análogamente, el quark $d$ es un
triplete
\[
  \begin{pmatrix}
    d_1\\
    d_2\\
    d_3
  \end{pmatrix}
\]
con carga $q=-\frac13e$.

El lagrangiano es
\[
  \mathcal L=
  i\bar u_a\not\!D u_a-m_u\bar u_a u_a
  +i\bar d_a\not\!D d_a-m_d\bar d_a d_a
  -\frac14F_{\mu\nu}F^{\mu\nu}
  -\frac14\Tr(G_{\mu\nu}G^{\mu\nu}),
\]
con
\[
  D_\mu u^a
  =\partial_\mu u^a-\frac23ieA_\mu u^a
   -igG_\mu^{(i)}\lambda^{(i)a}{}_{b}u^b,
\]
\[
  D_\mu d^a
  =\partial_\mu d^a+\frac13ieA_\mu d^a
   -igG_\mu^{(i)}\lambda^{(i)a}{}_{b}d^b.
\]

Las matrices $\lambda$ son los generadores de $SU(3)$ en la representación
fundamental, en la cual los elementos de $SU(3)$ son matrices unitarias
$3\times3$ con determinante $1$:
\[
\begin{gathered}
  \lambda^{(1)}=
  \begin{pmatrix}0&1&0\\1&0&0\\0&0&0\end{pmatrix},
  \quad
  \lambda^{(2)}=
  \begin{pmatrix}0&-i&0\\i&0&0\\0&0&0\end{pmatrix},
  \quad
  \lambda^{(3)}=
  \begin{pmatrix}1&0&0\\0&-1&0\\0&0&0\end{pmatrix},
  \quad
  \lambda^{(4)}=
  \begin{pmatrix}0&0&1\\0&0&0\\1&0&0\end{pmatrix},\\[0.8em]
  \lambda^{(5)}=
  \begin{pmatrix}0&0&-i\\0&0&0\\i&0&0\end{pmatrix},
  \quad
  \lambda^{(6)}=
  \begin{pmatrix}0&0&0\\0&0&1\\0&1&0\end{pmatrix},
  \quad
  \lambda^{(7)}=
  \begin{pmatrix}0&0&0\\0&0&-i\\0&i&0\end{pmatrix},
  \quad
  \lambda^{(8)}=
  \frac1{\sqrt3}
  \begin{pmatrix}1&0&0\\0&1&0\\0&0&-2\end{pmatrix}.
\end{gathered}
\]

A los ocho campos de gauge $G_\mu^{(i)}$ se los llama \concept{gluones} y son
responsables de mantener unidos dos $u$ y un $d$ para formar un protón, o dos
$d$ y un $u$ para formar un neutrón, a pesar de la repulsión electromagnética.

No podemos ver esta interacción $SU(3)$ (\concept{cromodinámica}) a nivel
macroscópico debido a un fenómeno cuántico. Resulta que a mayores distancias la
constante $g$ crece. Esto hace que los quarks no estén libres y sólo formen
objetos escalares $SU(3)$.

\end{document}
